\documentclass[11pt]{amsart}

\usepackage[T1]{fontenc}
\usepackage{microtype}
\usepackage{amsmath,amssymb}
\usepackage[hidelinks]{hyperref}

\newtheorem{theorem}{Theorem}
\theoremstyle{remark}
\newtheorem{remark}[theorem]{Remark}

\newcommand{\R}{\mathbb{R}}
\newcommand{\Rp}{\mathbb{R}_{+}}
\newcommand{\conv}{\operatorname{conv}}
\newcommand{\area}{\operatorname{area}}

\title[A counterexample to Conjecture 3]
{A Counterexample to Conjecture 3 of Fradelizi--Manui--Meyer--Ndiaye}
\author{Owen Shen}
\address{Operations Research Center, Massachusetts Institute of Technology,
Cambridge, MA 02139, USA}
\email{owenshen@mit.edu}
\date{}

\subjclass[2020]{Primary 52A20; Secondary 52A40}
\keywords{unconditional convex body, orthogonal projection, positive orthant,
Rogers--Shephard inequality}

\hypersetup{
  pdftitle={A Counterexample to Conjecture 3 of Fradelizi--Manui--Meyer--Ndiaye},
  pdfauthor={Owen Shen},
  pdfsubject={A counterexample for projections of unconditional convex bodies},
  pdfkeywords={unconditional convex body, orthogonal projection, positive orthant, Rogers--Shephard inequality}
}

\begin{document}

\begin{abstract}
Fradelizi, Manui, Meyer, and Ndiaye conjectured that
\[
  |P_EK|_n\le 2^n\bigl|P_E(K\cap\R_+^m)\bigr|_n
\]
for every unconditional convex body $K\subset\R^m$ and every $n$-dimensional
subspace $E\subset\R^m$. We give an exact counterexample for $(m,n)=(4,2)$,
with ratio $1399/336>4$. A Cartesian-product construction yields
counterexamples for every $m\ge4$ and $2\le n\le m-2$. Together with the
positive cases established in the cited work and the elementary case $n=m$,
this determines exactly the dimension pairs for which the inequality holds
universally.
\end{abstract}

\maketitle

\section{The conjecture and result}

A convex body $K\subset\R^m$ is \emph{unconditional} if it is invariant under
all coordinate sign changes. Let $P_E$ denote orthogonal projection onto a
subspace $E$, and let $|\cdot|_n$ denote $n$-dimensional Lebesgue measure on
$E$. Fradelizi, Manui, Meyer, and Ndiaye posed the following inequality as
Conjecture~3 in \cite{FMMN}:
\begin{equation}\label{eq:conjecture}
  |P_EK|_n\le 2^n\bigl|P_E(K\cap\R_+^m)\bigr|_n,
  \qquad \R_+^m=[0,\infty)^m.
\end{equation}

\begin{theorem}\label{thm:main}
For every $m\ge4$ and every $2\le n\le m-2$, there are an unconditional convex
body $K\subset\R^m$ and an $n$-dimensional subspace $E\subset\R^m$ for which
\[
  |P_EK|_n>2^n\bigl|P_E(K\cap\R_+^m)\bigr|_n.
\]
Consequently, for $1\le n\le m$, inequality~\eqref{eq:conjecture} holds for
all such $K$ and $E$ if and only if $n\in\{1,m-1,m\}$. In particular,
Conjecture~3 of \cite{FMMN} is false.
\end{theorem}

\section{The counterexample}

Set
\begin{equation}\label{eq:K0}
  K_0=\left\{x\in\R^4:
  4|x_1|+3|x_3|+3|x_4|\le1,\quad
  3|x_2|+4|x_3|\le1\right\},
\end{equation}
and let
\[
  b_1=(3,-1,1,2),\qquad b_2=(-3,1,2,0),\qquad
  E_0=\operatorname{span}\{b_1,b_2\}.
\]
The set $K_0$ is a full-dimensional unconditional convex body. Let $B$ be the
$4\times2$ matrix with columns $b_1,b_2$, and define
\[
  \Phi:\R^4\longrightarrow\R^2,\qquad \Phi(x)=B^{\mathsf T}x.
\]
Since
\[
  G=B^{\mathsf T}B=
  \begin{pmatrix}15&-8\\-8&14\end{pmatrix},
  \qquad \det G=146,
\]
one has, for every compact convex set $A\subset\R^4$,
\begin{equation}\label{eq:chart}
  \area\bigl(\Phi(A)\bigr)=\sqrt{146}\,|P_{E_0}A|_2.
\end{equation}
Indeed, $\Phi=\Phi\circ P_{E_0}$; if $y=Bc\in E_0$, then the Euclidean area
element is $\sqrt{\det G}\,dc$, whereas $\Phi(y)=Gc$ has area element
$(\det G)\,dc$.

Write $K_+=K_0\cap\R_+^4$. The six inequalities defining $K_+$ have fifteen
four-element active sets. Nine give feasible nonsingular intersections,
namely
\begin{align*}
V_+=\bigl\{&
(0,0,0,0),
(0,0,0,\tfrac13),
(0,0,\tfrac14,0),
(0,0,\tfrac14,\tfrac1{12}),
(0,\tfrac13,0,0),\\
&
(0,\tfrac13,0,\tfrac13),
(\tfrac1{16},0,\tfrac14,0),
(\tfrac14,0,0,0),
(\tfrac14,\tfrac13,0,0)
\bigr\}.
\end{align*}
The other six choices consist of four singular systems and the infeasible
intersections $(0,0,\tfrac13,0)$ and $(0,-\tfrac19,\tfrac13,0)$. Thus $V_+$ is
the complete vertex set of $K_+$. Unconditionality gives
\begin{equation}\label{eq:signhull}
  K_0=\conv\{\varepsilon\odot v:
  v\in V_+,\ \varepsilon\in\{-1,1\}^4\},
\end{equation}
where $\odot$ denotes coordinatewise multiplication.

Set $Q_+=\Phi(K_+)$ and $Q=\Phi(K_0)$. Applying $\Phi$ to the generators above
gives the exact vertex descriptions
\begin{align}
Q_+=\conv\bigl\{&
(-\tfrac13,\tfrac13),
(\tfrac34,-\tfrac34),
(\tfrac23,0),
(\tfrac5{12},\tfrac12),
(\tfrac14,\tfrac12)
\bigr\},\label{eq:Qplus}\\[2mm]
Q=\conv\bigl\{&
(-\tfrac{13}{12},\tfrac{13}{12}),
(-1,\tfrac13),
(-\tfrac5{12},-\tfrac12),
(-\tfrac1{16},-\tfrac{11}{16}),\notag\\
&
(\tfrac{13}{12},-\tfrac{13}{12}),
(1,-\tfrac13),
(\tfrac5{12},\tfrac12),
(\tfrac1{16},\tfrac{11}{16})
\bigr\}.\label{eq:Q}
\end{align}
The points are listed counterclockwise. For a compact exact certificate of
these hulls, the corresponding half-space descriptions are
\begin{align}
Q_+=\{(s,t):\;&s+t\ge0,\quad 9s+t\le6,\quad 6s+3t\le4,\notag\\
&2t\le1,\quad -2s+7t\le3\},\label{eq:HQplus}
\end{align}
and
\begin{equation}\label{eq:HQ}
\begin{aligned}
Q=\{(s,t):\;&|27s+3t|\le26,\quad |30s+21t|\le23,\\
              &|36s+68t|\le49,\quad |19s+55t|\le39\}.
\end{aligned}
\end{equation}
Indeed, every projected generator satisfies the relevant inequalities, while
the vertices of the two half-space intersections are exactly the points in
\eqref{eq:Qplus} and \eqref{eq:Q}. Thus the hull calculations require only
exact rational arithmetic.

The shoelace formula gives
\[
  \area(Q_+)=\frac7{12},
  \qquad
  \area(Q)=\frac{1399}{576}.
\]
Combining this with \eqref{eq:chart},
\begin{equation}\label{eq:ratio}
  \frac{|P_{E_0}K_0|_2}{|P_{E_0}K_+|_2}
  =\frac{1399/576}{7/12}
  =\frac{1399}{336}
  =4+\frac{55}{336}>4.
\end{equation}
Hence $(K_0,E_0)$ is a counterexample to \eqref{eq:conjecture} for
$(m,n)=(4,2)$.

\section{Product extension}

\begin{proof}[Proof of Theorem~\ref{thm:main}]
Fix $m\ge4$ and $2\le n\le m-2$, and use the orthogonal coordinate
decomposition
\[
  \R^m=\R^4\oplus\R^{n-2}\oplus\R^{m-n-2},
\]
with zero-dimensional factors omitted. Define
\[
  K=K_0\times[-1,1]^{m-4},
  \qquad
  E=E_0\oplus\R^{n-2}\oplus\{0\}^{m-n-2}.
\]
Then $K$ is unconditional and $\dim E=n$. Moreover,
\begin{align*}
  P_EK&=P_{E_0}K_0\times[-1,1]^{n-2},\\
  P_E(K\cap\R_+^m)&=P_{E_0}K_+\times[0,1]^{n-2}.
\end{align*}
Therefore, by \eqref{eq:ratio},
\[
  \frac{|P_EK|_n}{|P_E(K\cap\R_+^m)|_n}
  =2^{n-2}\frac{1399}{336}
  >2^{n-2}\cdot4
  =2^n.
\]
For the classification, the cases $n=m-1$ and $n=1$ are proved in
\cite[Propositions~3 and~4]{FMMN}, respectively. If $n=m$, then $E=\R^m$ and
unconditionality partitions $K$, up to null boundaries, into $2^m$ congruent
orthant pieces, so \eqref{eq:conjecture} is an equality. These are all the
remaining cases.
\end{proof}

\medskip
\noindent\textit{Scope.}
The body $K_0$ is neither an $\ell_q$-ball nor invariant under coordinate
permutations. Thus the example disproves Conjecture~3 as stated, but it does
not address Conjecture~2 of \cite{FMMN} or the restriction of Conjecture~3 to
$1$-unconditional bodies in the terminology of that paper.

\begin{thebibliography}{1}

\bibitem{FMMN}
M.~Fradelizi, A.~Manui, M.~Meyer, and C.~S.~Ndiaye,
\emph{$L_p$-Rogers--Shephard type inequalities for $L_p$-zonoids and symmetric
bodies}, arXiv:2607.03582v1 [math.MG], 2026.

\end{thebibliography}

\end{document}